\subsection{Unitary representations of the group R}

In this section, E denotes a complex Hilbert space. The group R is equipped with Lebesgue measure.

Lemma 11. --- Let u be a partial self-adjoint operator on E. For \(t \in \mathbf{R}\) set \(\varrho(t) = e^{itu} \in \mathcal{L}(\mathrm{E})\) . Then the map \(\varrho\) is a unitary representation of \(\mathbf{R}\) .

The operator \(\varrho(t)\) is defined by the universally measurable functional calculus (def. 5 of IV, p. 272); it is an endomorphism of E since the function \(x \mapsto e^{itx}\) is bounded on \(\mathrm{Sp}(u)\) (prop. 5, \(a\) ) of IV, p. 275).

According to prop. 5 of IV, p. 275, we have \(\varrho(0) = 1_{\mathrm{E}}\) , \(\varrho(t)^* = \varrho(-t)\) for all \(t \in \mathbf{R}\) and \(\varrho(t_1 + t_2) = \varrho(t_1)\varrho(t_2)\) for all \(t_1\) and \(t_2\) in \(\mathbf{R}\) . In particular, the endomorphism \(\varrho(t)\) is unitary for every \(t \in \mathbf{R}\) . For every \(x \in \mathrm{E}\) , the map from \(\mathbf{R}\) to E defined by \(x \mapsto \varrho(t)x\) is continuous at \(t = 0\) by prop. 6 of IV, p. 276, hence \(\varrho\) is a unitary representation of \(\mathbf{R}\) in E (V, p. 380, lemma 4).

Lemma 12. --- Let \(\varrho\) a unitary representation of \(\mathbf{R}\) in E. Let D be the set of elements \(x\) of E such that the map \(\psi_x: t \mapsto \varrho(t)x\) is differentiable at 0. The map \(u\) of D into E given by \(x \mapsto i^{-1} \psi_x'(0)\) defines a symmetric partial operator on E.

Let \(f \in \mathcal{D}(\mathbf{R})\) and \(x \in \mathrm{E}\) . Set \(y = \varrho(f)x\) . For every \(h \in \mathbf{R}\) , we have

\[
\begin{array}{r l} & {\psi_ {y} (h) = \varrho (h) \varrho (f) x = \varrho (f) \varrho (h) x} \\ & {\qquad = \int_ {\mathbf {R}} f (t) \varrho (t + h) x d t = \int_ {\mathbf {R}} f (t - h) \psi_ {x} (t) d t.} \end{array}
\]

The function of \(R^{2}\) to C defined by \((t,h)\mapsto f(t-h)\) is infinitely differentiable; its derivative with respect to h is the function defined by \((t,h)\mapsto -f'(t-h)\) , which is bounded by \(\|f'\|_{\infty}\) . When h=0, this derivative is zero for t outside a compact set. Since \(\|\psi_{x}(t)\|=\|x\|\) for all t, it follows from prop. 2 of IV, p. 197 that the function \(\psi_{y}\) is differentiable at 0 and that its derivative is given by

\[
\psi_ {y} ^ {\prime} (0) = - \int_ {\mathbf {R}} f ^ {\prime} (t) \psi_ {x} (t) d t = - \varrho (f ^ {\prime}) x.
\]

We therefore have \(y \in D\). Since the space \(\mathcal{D}(\mathbf{R})\) is dense in \(\mathrm{L}^{1}(\mathbf{R})\), it follows from INT, VIII, p. 139, § 2, no. 7, prop. 10, together with (prop. 4 of IV, p. 202), that the space D is dense in E.

Let \(x_{1}\) and \(x_{2}\) in D. We compute

\[
\begin{array}{c} \langle u (x _ {1}) \mid x _ {2} \rangle = \lim _ {h \to 0} \frac {i}{h} \langle (\varrho (h) - 1 _ {\mathrm{E}}) x _ {1} \mid x _ {2} \rangle \\ = \lim _ {h \to 0} \frac {i}{h} \langle x _ {1} \mid (\varrho (- h) - 1 _ {\mathrm{E}}) x _ {2} \rangle = \langle x _ {1} \mid u (x _ {2}) \rangle . \end{array}
\]

Therefore, the partial operator \(u\) is symmetric.

DEFINITION 5. --- Let \(\varrho\) a unitary representation of \(\mathbf{R}\) in E. The symmetric partial operator defined in Lemma 12 is called the infinitesimal generator of \(\varrho\) .

THEOREM 1 (Stone). --- The mapping σ which to a unitary representation ρ of R in E associates its infinitesimal generator u defines a bijection from the set of unitary representations of R in E onto the set of partial self-adjoint operators on E. The inverse bijection τ associates to a partial self-adjoint operator the unitary representation \(t \mapsto e^{itu}\) .

Let us first prove a few lemmas.

Lemma 13. --- Let \(\varrho\) a unitary representation of \(\mathbf{R}\) in E and let \(u\) its infinitesimal generator.

\begin{enumerate}
\def\labelenumi{\alph{enumi})}
\item
  The domain of u is the set of \(x \in E\) such that the mapping \(\psi_{x} \colon t \mapsto \varrho(t)x\) is differentiable on R;
\item
  For every \(t \in \mathbf{R}\) and every \(x \in \mathrm{dom}(u)\) , we have \(\varrho(t)x \in \mathrm{dom}(u)\) and \(\psi_x'(t) = i u(\varrho(t)x)\) ;
\item
  The partial operator u is essentially self-adjoint.
\end{enumerate}

For every \(x \in \mathrm{E}\) , denote \(\psi_x\) the map from \(\mathbf{R}\) in \(\mathrm{E}\) defined by \(\psi_x(t) = \varrho(t)x\) . For every \(t \in \mathbf{R}\) and every \(h \in \mathbf{R}\) , we have

\[
\psi_ {x} (t + h) - \psi_ {x} (t) = \varrho (t) \big (\psi_ {x} (h) - \psi_ {x} (0) \big),
\]

which proves that \(\operatorname{dom}(u)\) is the space of elements \(x \in E\) such that \(\psi_{x}\) is differentiable on R and establishes that \(\psi_{x}^{\prime}(t) = \varrho(t)\psi_{x}^{\prime}(0) = i\varrho(t)(u(x))\) for all \(t \in R\) .

Let \(x \in E\) and \(t \in R\) . We have \(\psi_{\varrho(t)x}(s) = \psi_x(s + t)\) for all \(s \in R\) . Consequently, one has \(\varrho(t)x \in \text{dom}(u)\) if \(x \in \text{dom}(u)\) , and moreover \(u(\varrho(t)x) = i^{-1}\psi_x'(t) = \varrho(t)u(x)\) by a). We then obtain assertion b).

Let us prove c). Let \(x \in \text{Ker}(u^{*} - i1_{\text{E}})\) . Let us prove that x = 0. Let \(y \in \text{dom}(u)\) , and let f be the function on R defined by \(f(t) = \langle \psi_{y}(t) | x \rangle\) for \(t \in R\) . The function f is bounded since \(\| \psi_{y}(t) \| = \| y \|\) for all \(t \in R\) ; it is differentiable on R and, for every \(t \in R\) , assertion b) implies

\[
\begin{array}{r l} f ^ {\prime} (t) = \langle \psi_ {y} ^ {\prime} (t) | x \rangle & = \langle i u (\varrho (t) y) | x \rangle \\ & = - i \langle \varrho (t) y | u ^ {*} (x) \rangle = - i \langle \psi_ {y} (t) | i x \rangle = f (t) \end{array}
\]

since \(u^{*}(x)=ix\) . We therefore have \(f(t)=f(0)e^{t}\) for all \(t\in R\) (FVR, IV, p. 27). Since f is bounded, the function f is zero and, in particular, we have \(\langle y|x\rangle=f(0)=0\) . Since the space dom(u) is dense in E, we conclude that x=0. Similarly, we show that \(\operatorname{Ker}(u^{*}+i1_{\mathrm{E}})\) is reduced to 0. By Cor. 3 of IV, p. 261, the partial operator u is therefore essentially self-adjoint.

Lemma 14. --- Let u be a partial self-adjoint operator on E, and let \(\varrho(t) = e^{itu}\) the unitary representation of R defined by u. The infinitesimal generator of \(\varrho\) is equal to u.

Let \(x \in \text{dom}(u)\) . Set \(\psi_{x}(t) = \varrho(t)x\) for all \(t \in R\) . For every nonzero real number h, we have

\[
\frac {1}{h} \left(\psi_ {x} (t + h) - \psi_ {x} (t)\right) = \left(\frac {1}{h} \left(e ^ {i h u} - 1 _ {\mathrm{E}}\right)\right) e ^ {i t u} x = \left(\frac {1}{h} \left(e ^ {i h u} - 1 _ {\mathrm{E}}\right)\right) \varrho (t) x.
\]

As h tends to 0, we have

\[
\frac {1}{h} (e ^ {i h t} - 1) \rightarrow i t
\]

for all \(t \in R\) . Moreover,

\[
\left| \frac {1}{h} (e ^ {i h t} - 1) \right| = | t | \left| \frac {1}{h} \int_ {0} ^ {h} e ^ {i t s} d s \right| \leqslant | t |.
\]

It then follows from prop. 6 of IV, p. 276 that the function \(\psi_x\) is differentiable on \(\mathbf{R}\) and satisfies \(\psi_x'(t) = i u(\varrho(t)x)\) for all \(t \in \mathbf{R}\) . Consequently, the domain of \(u\) is contained in the set of \(x \in E\) such that \(\psi_x\) is differentiable at 0 and that we then have \(\psi_x'(0) = i u(x)\) . This means by definition that the infinitesimal generator of \(\varrho\) is an extension of \(u\) . These two operators are therefore equal since they are symmetric and \(u\) is self-adjoint (IV, p. 238, remark 5).

Lemma 15. --- Let \(\varrho\) a unitary representation of \(\mathbf{R}\) in E. Then the infinitesimal generator \(u\) of \(\varrho\) is self-adjoint and \(\varrho(t) = e^{itu}\) for all \(t \in \mathbf{R}\) .

By Lemma 13(c), the partial operator u is essentially self-adjoint. Its closure \(\overline{u}\) is therefore a self-adjoint operator. Denote \(\pi\) the unitary representation of R defined by \(\pi(t) = e^{it\overline{u}}\) (lemma 11).

For every \(x \in \mathrm{E}\) , denote \(\psi_x\) (resp. \(\widetilde{\psi}_x\) ) the map from \(\mathbf{R}\) in \(\mathrm{E}\) defined by \(\psi_x(t) = \varrho(t)x\) (respectively by \(\widetilde{\psi}_x(t) = \pi(t)x\) ).

By Lemma 13, parts (a) and (b), the space \(\mathrm{dom}(u)\) is the subspace of E consisting of the elements \(x \in \mathrm{E}\) such that the mapping \(\psi_x\) is differentiable on \(\mathbf{R}\) , and for all \(x \in \mathrm{dom}(u)\) and every \(t \in \mathbf{R}\) , we have \(\psi_x'(t) = i u (\psi_x(t))\) . Likewise, the space \(\mathrm{dom}(\overline{u})\) is the subspace of E consisting of the elements \(x \in \mathrm{E}\) such that the mapping \(\widetilde{\psi}_x\) is differentiable on \(\mathbf{R}\) , and for all \(x \in \mathrm{dom}(\overline{u})\) and every \(t \in \mathbf{R}\) , we have \(\widetilde{\psi}_x'(t) = i \overline{u} (\widetilde{\psi}_x(t))\) .

Let \(x \in \text{dom}(u) \subset \text{dom}(\overline{u})\) . Set \(f = \psi_x - \widetilde{\psi}_x\) . This is a differentiable function from \(\mathbf{R}\) into E. For every \(t \in \mathbf{R}\) , it follows that

\[
f ^ {\prime} (t) = i u (\psi_ {x} (t)) - i \overline {{u}} (\widetilde {\psi} _ {x} (t)) = i \overline {{u}} (f (t))
\]

since \(\psi_x(t) \in \mathrm{dom}(u)\) and \(u \subset \overline{u}\) . Set \(g = \|f\|^2\) ; it is a differentiable map from \(\mathbf{R}\) in \(\mathbf{R}\) such that \(g(0) = 0\) . For \(t \in \mathbf{R}\) , we obtain by FVR, I, p. 28, prop. 2

\[
\begin{array}{r l} g ^ {\prime} (t) & = \langle f ^ {\prime} (t) | f (t) \rangle + \langle f (t) | f ^ {\prime} (t) \rangle \\ & = \langle i \overline {{u}} (f (t)) | f (t) \rangle + \langle f (t) | i \overline {{u}} (f (t)) \rangle = 0 \end{array}
\]

since \(\overline{u}\) is self-adjoint. Hence f = 0, so \(\varrho(t)x = \pi(t)x\) for all \(t \in R\) . The continuous endomorphisms \(\varrho(t)\) and \(\pi(t)\) of E coincide on dom(u), and are therefore equal for all \(t \in R\) . Thus, we have \(\pi = \varrho\) ; since \(\overline{u}\) is the infinitesimal generator of \(\pi\) (Lemma 14), we have \(\overline{u} = u\) , which proves that u is self-adjoint.

We can now prove Theorem 1.

The maps \(\sigma\) and \(\tau\) are well-defined (Lemma 15 and Lemma 11, respectively).

Let \(\varrho\) a unitary representation of \(\mathbf{R}\) . Denote by \(u\) its infinitesimal generator. The relation \(\varrho(t) = e^{itu}\) for all \(t \in \mathbf{R}\) (lemma 15) proves that \(\tau \circ \sigma\) is the identity map.

Let u be a partial self-adjoint operator on E. Lemma 14 shows that the infinitesimal generator of the unitary representation \(t \mapsto e^{itu}\) is equal to u, hence \(\sigma \circ \tau\) is the identity map.

Let u be a partial self-adjoint operator on E and let \(\varrho(t) = e^{itu}\) for \(t \in R\) the unitary representation of R in E associated with it.

Let \(x \in \text{dom}(u)\) . The equation \(\partial_{t} \varrho(t)x = i u(\varrho(t)x)\) which is then satisfied (Lemma 13, b)) is called the Schrödinger equation.

COROLLARY. --- Let \(\varrho\) a unitary representation of \(\mathbf{R}\) in a Hilbert space E. There exists a locally compact space X, a positive measure \(\mu\) on X and a continuous function \(g\) on X with real values such that \(\varrho\) is isomorphic to the representation \(\pi\) of \(\mathbf{R}\) in \(\mathrm{L}^2 (\mathrm{X},\mu)\) defined by \(\pi (t)f = e^{itg}f\) for all \(t\in \mathbf{R}\) and every \(f\in \mathrm{L}^2 (\mathrm{X},\mu)\) .

Let \(u\) the self-adjoint operator on E such that \(\varrho(t) = e^{itu}\) for all \(t \in \mathbf{R}\) (Theorem 1). There exists a locally compact space X, a positive measure \(\mu\) on X, an isometric isomorphism \(\theta\) of \(L^2(X, \mu)\) on E and a continuous function \(g\) on X with real values such that \(u = \theta \circ m_g \circ \theta^{-1}\) (theorem 1 of IV, p. 266). The assertion follows from the formula \(e^{itu} = \theta \circ e^{itm_g} \circ \theta^{-1} = \theta \circ m_{e^{itg}} \circ \theta^{-1}\) (lemma 4 of IV, p. 269).

Remark. --- Suppose that u is an endomorphism of E. The unitary representation of R in E defined by \(\varrho(t) = e^{itu}\) then verifies the inequality \(\|\varrho(t) - 1_{\mathrm{E}}\| \leqslant |t||u\|\) for all \(t \in R\) , and the map \(\varrho\) from R into the Banach space \(\mathcal{L}(\mathrm{E})\) is the unique solution of the linear differential equation

\[
\frac {1}{i} \frac {d \varrho}{d t} = u \circ \varrho
\]

(cf. FVR, IV, p. 26, §6).

Example. --- Let \(\varrho\) the regular representation of \(\mathbf{R}\) in \(\mathrm{L}^2 (\mathbf{R})\) . The infinitesimal generator of \(\varrho\) is the closure of the differential operator with domain \(\mathcal{D}(\mathbf{R})\) defined by \(f\mapsto -if'\) .

